\section{Experimental Design}\label{sec:design}
The comparison is paired. For every dataset and rule, the input matrix, the gene set, the cap
on the candidate range and the scoring are fixed, and the method and its grid differ.

\subsection{Data}
We use the ten labelled datasets of the scMINER benchmark in the processed form the
authors released, so differences from their reported numbers cannot come from our
preprocessing. They span 124 to 13,605 cells and 3 to 11 cell types, the range over which
the published claim was made (Table~\ref{tab:data}). PBMC14K's labels are cell types fixed by
antibody sorting before sequencing \cite{zheng2017massively}. Several small datasets are
labelled by experimental condition, such as cell-cycle phase on Buettner, culture medium on
Kolodziejczyk and time point on Klein, while Zeisel's labels came from the original authors
clustering the same data \cite{zeisel2015cell}, which flatters any method resembling theirs.

\subsection{What is held fixed}
Both methods receive the same matrix with all genes retained, and both choose among
settings whose full-data partition has 2 to 20 clusters. That caps the range without matching
it. MICA's $k$ covers all of it on every dataset, but the baseline's resolution grid reaches
only part of it on most, so we also re-apply the rule to MICA's saved scores restricted to the
counts the baseline produced (Table~\ref{tab:common}). That restriction is made after the
results were seen and only to MICA, so it is a check on what the result depends on, not a
second experiment. The true number of types is never supplied. Both see the same $B = 25$
subsamples within a seed, which makes the comparison paired at the level of the draw.

chooseR specifies $B = 100$. We use 25 because the rule costs $B + 1$ clusterings per
candidate setting, so 100 would take roughly four times the compute for every result here.
The reduction widens the bootstrap intervals and lowers the threshold, which could change
which settings clear it, so rather than assume it is harmless we test the selection's
sensitivity to the threshold directly in Section~\ref{sec:results}.

The protocol was fixed before any partition was scored: the candidate grids, 25 subsamples
of 80 percent at three seeds, and the selection step. chooseR's own reading is the primary
rule. The plain maximum, specified with it and computed from the same scores in the same run,
is a second rule, and comparing the two measures how much the result depends on the
tie-break.

\subsection{Use of labels}
Labels are used in three places, none of which feeds a label-free choice. They score the
chosen partition, they fix the setting in the reproduction of the published protocol, an
upper bound rather than a result a user could obtain, and they pick the resolution in the
graph comparison so that only the graph differs. The transfer rule in our plan picks each
dataset's setting by its ARI on the others, so it is not label-free and will be reported
apart.

\subsection{Runs and variation}
Every stochastic comparison is repeated. The stability rule is run with three independent
seeds, which redraw the subsamples and reinitialise the clustering, and we report the
median across those runs with the full range of chosen cluster counts. Per-dataset ARI
ranges are in \texttt{results/a3\_seed\_variation.csv}. MICA's $k$-means uses 50 restarts
within a run, enough to reach its own optimum reliably, which are restarts rather than
repeats. The default rule was run at three seeds and the silhouette rule at one.
Where a range is reported it is the observed minimum and maximum, not a confidence
interval, because three runs are too few to justify one.

\subsection{Measures}
Agreement with the labels is the ARI in the main table, and the adjusted mutual information
for every run is in Table~\ref{tab:ami}. The ARI is the stricter of the two here in both
directions: across MICA's 54 stability picks the AMI sits above it by 0.17 on average when the
rule chose too many clusters and by 0.12 when it chose too few, which it did in 35 of them. We also report $|P_{\hat{s}}|$, the number of clusters the rule chose, because a
rule can reach a respectable ARI with a badly wrong number of clusters, and a user would
notice the latter.

\subsection{Cost}
We report the cost of selecting a setting apart from the cost of clustering once, since a
user pays both. Cost is wall-clock time on one workstation with 120 cores and 754 GB, every
job on 8 threads and no GPU. Selection counts the candidate partitions, the subsample reruns
and the scoring. MICA's one-off distance matrix and the number of clustering calls are given
separately.

\subsection{Graph comparison}
The faster neighbour graph is compared with MICA's own on PBMC14K and on 100,000 cells of the
Allen brain atlas, both fed to MICA's own node2vec and Louvain so that only the graph differs.
The candidate count follows one declared rule: the smallest value on a doubling grid whose
recall on 500 brute-forced query cells reaches that of MICA's search. It gives the counts used
in Assignment 2, 800 and 3,200, and stays fixed for any further dataset. Recall here is
agreement with exact neighbours, so setting it uses no labels.
